\documentclass[12pt]{article}
\usepackage[utf8]{inputenc}
\usepackage{geometry}
\geometry{a4paper, margin=1in}
\usepackage{booktabs}
\usepackage{enumitem}
\usepackage{titlesec}
\usepackage{amsmath}
\usepackage{hyperref}

\titleformat{\section}{\large\bfseries}{\thesection.}{0.5em}{}
\titleformat{\subsection}{\normalsize\bfseries}{\thesubsection}{0.5em}{}

\title{\textbf{Project Proposal: AI Racing Challenge} \\ \large Human vs. AI Racing at the Limits of Handling: Pure Pursuit, MPCC, and Reinforcement Learning}
\author{Team Proposal}
\date{\today}

\begin{document}

\maketitle

\section{Project Overview}
This project aims to develop a simulated 2D autonomous racing environment to benchmark three generations of vehicle control architectures operating at physical handling limits. Operating within a high-frequency simulation loop allows the team to exploit perfect state information to evaluate vehicle dynamics at the threshold of friction saturation.

The simulation will pit three paradigms against each other:
\begin{enumerate}[noitemsep]
    \item \textbf{Geometric Control (Pure Pursuit):} A classical tracking baseline using a curvature-based velocity profile as a reference point for non-optimization control.
    \item \textbf{Model Predictive Contouring Control (MPCC):} An optimal control loop (Liniger et al., 2015) that maximizes track progress along a contouring horizon while adhering to non-linear dynamic constraints.
    \item \textbf{Reinforcement Learning (Proximal Policy Optimization):} A model-free data-driven agent that learns competitive racing lines and input trajectories through continuous training.
\end{enumerate}
Performance will be evaluated primarily on lap time, supported by telemetry analysis of racing line choices, input smoothness, and track generalization.

\section{Detailed Simulation Framework \& Plant Dynamics}
The project will be implemented in Python using Pygame for rendering and custom numerical integrations for physics updates.

\subsection*{Dynamic Bicycle Model with Pacejka Non-Linear Tires}
To capture realistic vehicle handling near grip boundaries, we implement a 2-DOF Dynamic Bicycle Model with non-linear Pacejka tire dynamics. The full state vector is defined as:
\[ x_k = [x, y, \psi, v_x, v_y, r]^T \]
where $(x, y)$ are global coordinates, $\psi$ is heading angle, $v_x$ is longitudinal velocity, $v_y$ is lateral velocity, and $r$ is yaw rate ($\dot{\psi}$). Controls are acceleration/braking $a_x$ and front steering angle $\delta$.

\textbf{Tire Saturation Model:} Front and rear slip angles ($\alpha_f, \alpha_r$) are calculated as:
\[ \alpha_f = \delta - \arctan\left(\frac{v_y + l_f r}{v_x}\right), \quad \alpha_r = -\arctan\left(\frac{v_y - l_r r}{v_x}\right) \]
To model tire force saturation during aggressive cornering, lateral forces ($F_{yf}, F_{yr}$) are computed using Pacejka's Magic Formula:
\[ F_{y} = D \sin\left(C \arctan\left(B \alpha - E (B \alpha - \arctan(B \alpha))\right)\right) \]
where $B, C, D, E$ represent empirical shape, stiffness, peak, and curvature factors.

\textbf{Low-Speed Kinematic Switching:} To prevent numerical division-by-zero singularities when $v_x \approx 0$, the update equations switch to a kinematic bicycle model at low speeds ($v_x < v_{\text{min}}$):
\[ \dot{x} = v_x \cos\psi, \quad \dot{y} = v_x \sin\psi, \quad \dot{\psi} = \frac{v_x}{L} \tan\delta \]
Above $v_{\text{min}}$, full dynamic updates govern motion:
\begin{align*}
v_{x, k+1} &= v_{x, k} + (a_x + v_y r) \Delta t \\
v_{y, k+1} &= v_{y, k} + \left( \frac{F_{yf}\cos\delta + F_{yr}}{m} - v_x r \right) \Delta t \\
r_{k+1} &= r_k + \left( \frac{l_f F_{yf}\cos\delta - l_r F_{yr}}{I_z} \right) \Delta t
\end{align*}

\section{Controller Architectures}

\subsection*{1. Baseline: Pure Pursuit with Curvature-Based Velocity Profile}
Pure Pursuit serves as our baseline reference. It calculates steering commands relative to a look-ahead point distance $L_t$:
\[ \delta = \arctan\left(\frac{2L \cdot x_{\text{local}}}{L_t^2}\right) \]
To remain functional on a race track, it is paired with a pre-computed curvature velocity profile $v_{\text{ref}}(s) = \sqrt{\mu g / \kappa(s)}$, where $\kappa(s)$ is local path curvature and $\mu$ is friction coefficient.

\subsection*{2. Optimal Control: Model Predictive Contouring Control (MPCC)}
Following Liniger et al. (2015), the MPC optimization shifts from trajectory-tracking to Model Predictive Contouring Control (MPCC). Rather than minimizing lateral error relative to a fixed time-indexed trajectory, MPCC maximizes progress along the track centerline arc-length $s$ while penalizing deviation within track boundaries.

The objective function $J$ over dynamic prediction horizon $N$ is:
\[ J = \sum_{t=0}^{N} \left( -w_s \Delta s(t) + w_l e_l(t)^2 + w_c e_c(t)^2 + w_u \|\Delta u(t)\|^2_2 \right) \]
where $\Delta s(t)$ is forward progress along the track, $e_l(t)$ and $e_c(t)$ are lag and contouring errors relative to the reference path, and $\Delta u(t) = [\Delta\delta, \Delta a_x]^T$ penalizes abrupt actuator shifts.

\textbf{Track Boundary Constraints:} Hard constraints enforce track boundaries rather than soft cost penalties:
\[ |e_{ct}(t)| \le \frac{w}{2} - \frac{d_{\text{car}}}{2}, \quad \forall t \in [0, N] \]
where $w$ is track width and $d_{\text{car}}$ is vehicle width. Actuator constraints enforce mechanical limits ($|\delta| \le \delta_{\text{max}}, |a_x| \le a_{\text{max}}$). The optimization problem is solved using \textbf{CasADi} with the \textbf{IPOPT} solver.

\subsection*{3. Learning-Based Control: Proximal Policy Optimization (PPO)}
% 1. Setup (1 sentence)
-- The policy (actor) : a small neural network . observation goes in =; out comes a 

\begin{itemize}[noitemsep]
    \item \textbf{Observation:} % what it sees + why
    the 
    \item \textbf{Action:} % what it outputs + why

    \item \textbf{Reward and termination:} % reward, why progress, when episodes end

    \item \textbf{Training:} % curriculum, seeds
\end{itemize}

\section{Performance Evaluation Matrix}
Lap time is established as the primary benchmark. Cross-track error and actuator metrics serve to analyze line strategy and controller execution.

\begin{table}[h!]
\centering
\small
\begin{tabular}{lp{10cm}}
\toprule
\textbf{Metric} & \textbf{What We Plan to Measure} \\ \midrule
Lap Time (Primary) & Absolute time required to complete full laps at racing speeds. \\ \addlinespace
Line Selection Characteristics & Profile of cross-track error ($e_{ct}$) to analyze geometric racing lines. \\ \addlinespace
Actuator Smoothness & RMS rate of change ($\Delta\delta, \Delta a_x$) to quantify input chatter. \\ \addlinespace
Track Completion Rate & Percentage of successful laps completed without boundary collisions. \\ \addlinespace
Zero-Velocity Stability & System stability transitioning out of full stops ($v_x = 0$) using kinematic switching. \\ \bottomrule
\end{tabular}
\caption{Racing Evaluation Metrics}
\end{table}

\section{Team Responsibilities}
\begin{table}[h!]
\centering
\small
\begin{tabular}{llp{7.5cm}}
\toprule
\textbf{Member} & \textbf{Main Responsibility} & \textbf{Tasks} \\ \midrule
Member 1 & Environment \& Dynamics & Implement Milliken dynamics, Pacejka tire formulas, kinematic switching for low speeds, and track maps. \\ \addlinespace
Member 2 & Optimal Control (MPCC) & Formulate the MPCC optimization problem in CasADi/IPOPT, code progress parameters, and set track limits. \\ \addlinespace
Member 3 & Reinforcement Learning & Build Gymnasium wrapper, engineer progress-focused reward structures, and train PPO agents. \\ \addlinespace
Member 4 & Integration \& Analytics & Implement Pure Pursuit baseline, integrate controller profiles, collect telemetry data, and create comparative plots. \\ \bottomrule
\end{tabular}
\caption{Team Task Allocations}
\end{table}

\section{Expected Project Timeline}
\begin{description}[leftmargin=3.5cm, style=nextline]
    \item[November 5] Finalize team assignments and outline architectural changes.
    \item[November 12] Present MPCC scope, Pacejka tire dynamics, and project plan to the professor.
    \item[November 19] Complete 2D simulation loop, kinematic switching, and initial track boundary implementations.
    \item[Nov 20--26] Complete Pure Pursuit baseline. Formulate CasADi MPCC optimization script and set up parallel PPO training.
    \item[Nov 27--Dec 3] Tune MPCC constraints and PPO progress reward weights. Verify tire saturation handling.
    \item[December 4--9] Merge controllers into comparative multi-car racing modes and log telemetry datasets.
    \item[December 10] Present racing simulation results and comparative lap time analyses.
    \item[December 11--17] Finalize analysis and submit the final project proposal report.
\end{description}

\end{document}
